\chapter{Screenshot Manifest}
\label{app:screenshots}

Every figure in this book that requires a capture from a running tool. Entries
disappear from the list at the end of this appendix automatically once the
corresponding file exists in \code{figures/screenshots/}.

% ============================================================
\section{How the mechanism works}

In the source, a pending capture is written as:

\begin{yamlcode}[caption={Requesting a screenshot}, label={lst:needscreenshot}]
\needscreenshot{vs-about-dialog}%
{Caption printed under the figure}%
{Precise description of what to capture.}
\end{yamlcode}

While \code{figures/screenshots/vs-about-dialog.png} is absent, a dashed
placeholder is printed in place carrying the capture instructions, and the request
is added to the list below. Drop the file in with the matching name and it is
included automatically --- no edit to the chapter is required.

To see what is outstanding without building the PDF:

\begin{shellcmd}
make shots
\end{shellcmd}

% ============================================================
\section{Capture conventions}

\begin{itemize}[leftmargin=1.4em]
\item PNG, light theme, at least 1600\,px wide.
\item Crop tightly to the relevant dialog or pane. A full-desktop capture scaled
      to a 17\,cm page is illegible, and illegible screenshots are worse than no
      screenshot.
\item Editor font at 100\% or larger. Read the capture at print size before
      accepting it.
\item Where a pane relationship is the point --- editor plus Solution Explorer,
      say --- include both, but nothing else.
\item Consistent window chrome and theme across the whole book. Decide once.
\item Filename must match the key exactly, including hyphens.
\end{itemize}

% ============================================================
\section{Redaction checklist}

Run through this before accepting any capture. Screenshots are the most common
route by which private information reaches a published technical book, because
the author is looking at the feature rather than the frame.

\begin{itemize}[leftmargin=1.4em]
\item No tenant names, subscription IDs, resource group names, or endpoint hosts
      that identify a real deployment.
\item No API keys, tokens, or connection strings --- including in a partially
      scrolled \code{secrets.json} or an environment variable pane.
\item No personal file paths. \code{C:\textbackslash Users\textbackslash yourname\textbackslash{}}
      appears in more title bars than you expect.
\item No email addresses, account names, or profile photographs in the IDE
      account area. Sign out, or crop it.
\item No unrelated open documents, tabs, or notifications.
\item No customer or employer identifiers in recent-files lists, window titles,
      or Git branch names.
\end{itemize}

\begin{warning}
Redact by \emph{recapturing}, not by drawing a black box over the image. Black
boxes advertise that something was there, and a surprising number of them are
applied at a layer that leaves the original pixels in the file.
\end{warning}

% ============================================================
\section{Working practice}

Two habits worth adopting early.

\textbf{Capture as you go, not at the end.} The screenshots in this book are
taken from a solution being built chapter by chapter. Recreating the exact state
of a Chapter~\ref{ch:memory} debugging session three months later, in order to
photograph it, is an afternoon you will resent.

\textbf{Recapture when the version changes.} A screenshot of a dialog from an IDE
two versions old is a stronger signal that the book is stale than any prose
disclaimer. Figure~\ref{fig:vs-about-dialog} in Chapter~\ref{ch:landscape} exists
precisely so a reader can tell which build everything else was taken against ---
which means it must be recaptured first when you upgrade.

% ============================================================
\section{Outstanding captures}

\listofshots
